% ============================================================
%  Forschungsprojekt I – Flag Football GPS-Analyse
%  Bericht im Format IEEE Transactions (Vorlage: IEEEtran)
%  Kompilieren: pdflatex bericht.tex (zweimal, wegen Verweisen)
% ============================================================
\documentclass[journal]{IEEEtran}

\usepackage[utf8]{inputenc}
\usepackage[T1]{fontenc}
\usepackage[ngerman]{babel}
\usepackage{amsmath}
\usepackage{graphicx}
\usepackage{booktabs}
\usepackage{cite}
\usepackage{url}

\begin{document}

\title{Automatische Sprint- und Antritt-Erkennung im Amateur-Flag-Football mit 1-Hz-GPS-Daten einer Sportuhr}

\author{Clemens Näther%
\thanks{Forschungsprojekt I, HTW Dresden, Trainingswissenschaft \& Sportinformatik.}}

\markboth{Forschungsprojekt I, HTW Dresden, 2026}{}

\maketitle

% ------------------------------------------------------------
\begin{abstract}
% TODO: (ca. 150--250 Wörter)
% Aufbau: Problem -- Forschungsfrage -- Methode -- wichtigstes Ergebnis -- Fazit.
\end{abstract}

\begin{IEEEkeywords}

\end{IEEEkeywords}

% ------------------------------------------------------------
\section{Einleitung}

\IEEEPARstart{F}{lag} Football wächst als Breiten- und Wettkampfsport, wissenschaftliche
Belastungsanalysen gibt es bisher jedoch kaum. Im Profisport werden Laufbelastungen
mit GPS-Systemen erfasst, die für Amateurvereine zu teuer sind. Viele Spieler tragen
aber bereits eine GPS-Sportuhr.

% TODO: 2--3 Sätze zur Lücke: Warum reicht das, was es gibt, nicht aus?

Daraus ergibt sich die Forschungsfrage dieser Arbeit:

\begin{quote}
\textit{Lassen sich aus 1-Hz-GPS-Daten einer Consumer-Sportuhr im Amateur-Flag-Football
Sprints und explosive Antritte automatisch erkennen, sodass sie für Coaching-Feedback
nutzbar sind?}
\end{quote}

% TODO: Beitrag der Arbeit in 2--3 Stichpunkten, z.B. Pipeline, Schwellen, Replay-Tool.

% ------------------------------------------------------------
\section{Verwandte Arbeiten}

% TODO: Kurz halten (1--2 Absätze).
% - GPS-Metriken und individuelle Geschwindigkeitszonen im Teamsport [1]
% - Validität von GPS bei kurzen Sprints / niedriger Abtastrate [2]
% - Segmentierung von Bewegungen mit Wearables [3]
% - Tracking-Daten im American Football [4]

% ------------------------------------------------------------
\section{Methodik}

\subsection{Datengrundlage}
Erfasst wurden 14 Einheiten eines Spielers mit einer Garmin-Uhr (1\,Hz):
8 Einheiten aus einem Uni-Kurs, 3 Vereinstrainings und 3 Turnierspiele.

% TODO: Uhrmodell, Spielposition, Zeitraum ergänzen.

\subsection{Vorverarbeitung}
Geschwindigkeitswerte über 12\,m/s werden als Ausreißer verworfen und linear
interpoliert. Anschließend wird die Geschwindigkeit mit einem zentrierten
gleitenden Mittelwert über drei Samples geglättet:
\begin{equation}
  \hat{v}(t) = \frac{1}{3}\sum_{k=-1}^{1} v(t+k).
\end{equation}

\subsection{Sprint- und Antritt-Erkennung}
Alle Schwellen beziehen sich auf die persönliche Maximalgeschwindigkeit
$v_{\max}$ (99,99.~Perzentil über alle Einheiten) \cite{sfs_metrics}.
Ein Lauf liegt vor, wenn
\begin{equation}
  \hat{v}(t) \geq 0{,}6 \cdot v_{\max}
\end{equation}
für mindestens 2\,s gilt \cite{sfs_validity}. Antritte werden getrennt über die
Beschleunigung $a(t) \geq 2\,\text{m/s}^2$ erkannt.

% TODO: Später ggf. Tabelle der Geschwindigkeitszonen aus doku.md übernehmen.

% ------------------------------------------------------------
\section{Ergebnisse}

% TODO: Erst füllen, wenn die Zahlen feststehen. Vorschlag:
% - Tabelle: Sprints und Antritte je Kontext (Uni-Kurs / Training / Turnier)
% - Abbildung: Speed-Verlauf einer Beispiel-Session (plot_sprints)

% ------------------------------------------------------------
\section{Diskussion}

% TODO: Forschungsfrage beantworten, dann Limitationen:
% - 1\,Hz ist für sehr kurze Aktionen grenzwertig
% - Schwellen noch nicht per Video validiert
% - nur ein Proband

% ------------------------------------------------------------
\section{Fazit und Ausblick}

% TODO: 1 Absatz Fazit; Ausblick: Spielzug-Segmentierung, Routen-Klassifikation, Dashboard.

% ------------------------------------------------------------
\begin{thebibliography}{9}

\bibitem{sfs_metrics}
Science for Sport, ``GPS wearables: Metrics and application.''
[Online]. Verfügbar: \url{https://www.scienceforsport.com/gps-wearables-metrics-and-application/}

\bibitem{sfs_validity}
Science for Sport, ``GPS wearables: Validity and reliability.''
[Online]. Verfügbar: \url{https://www.scienceforsport.com/gps-wearables-validity-and-reliability/}

\bibitem{kim2018}
W.~Kim and M.~Kim, ``On-line detection and segmentation of sports motions using a
wearable sensor,'' \textit{Sensors}, vol.~18, no.~3, p.~913, 2018,
doi: 10.3390/s18030913.

\bibitem{yurko2025}
R.~Yurko \textit{et al.}, ``Fractional tackles: Leveraging player tracking data for
within-play tackling evaluation,'' 2025. [Online]. Verfügbar:
\url{https://www.ncbi.nlm.nih.gov/pmc/articles/PMC11739690/}

\end{thebibliography}

\end{document}
